\chapter{I problemi incontrati, e come sono stati risolti}
\label{ch:appendix}

Un documento di riferimento che mostrasse soltanto la versione
funzionante nasconderebbe la propria prova migliore. Ogni test di
questo programma è costruito per poter dire \emph{no}: la sua banda
di superamento è derivata dall'errore misurato dello schema stesso
anziché scelta a mano, e ogni suite impianta corruzioni deliberate
--- i \emph{rigettatori} --- che devono farla fallire. Un difetto
colto dalla suite non è quindi un imbarazzo ma una misura: la prova
che lo strumento ha i denti, e un resoconto di che cosa fa la fisica
quando la si sbaglia. I capitoli principali mostrano ciascun passo
così come funziona oggi; questa appendice raccoglie, passo per passo,
che cosa è andato storto lungo la strada, come è stato diagnosticato
e che cosa è stato cambiato. Nulla qui è aneddoto --- ogni voce è
stata colta da un test e chiusa da una correzione misurata.

\section{Passo 1: la superficie di controllo non chiusa}
\label{app:sliver}

La spinta è calcolata come bilancio su una \emph{superficie di
controllo}: una superficie chiusa immaginaria tracciata attorno al
flusso, sulla quale si sommano la quantità di moto che la attraversa
e la pressione che vi spinge. Una verifica di quella contabilità è il
\emph{test di gauge} --- si somma ovunque la stessa costante alla
pressione e nulla deve cambiare, perché una costante integra a zero
su una superficie genuinamente chiusa. Fallì alla primissima
esecuzione (Capitolo~\ref{ch:score}).

Causa: l'integrale di parete della via 2 partiva dal primo
\emph{punto} di parete della maglia, lasciando scoperta una scheggia
microscopica di parete fra il labbro della gola e quel punto. La
superficie non era chiusa, quindi una costante non integrava più a
zero (il cerchio rosso nella Figura~\ref{fig:tworoutes}). Correzione:
chiudere la superficie al labbro della gola. Il residuo di quantità
di moto $R_{\rm mom}$ --- il disaccordo fra le due vie indipendenti
di calcolo della spinta --- si è \emph{stretto} una volta che la
scheggia ha smesso di nascondersi al suo interno: il difetto ne stava
assorbendo una parte del bilancio meccanico.

\section{Passo 6: il rigettatore corretto dalla fisica}
\label{app:f2rejector}

Un rigettatore è utile solo se il caso corrotto che pone è uno che il
test potrebbe plausibilmente far fallire. Il primo controllo negativo
per il test del primo integrale della Sezione~\ref{sec:o33} --- la
verifica che una particolare combinazione $f_2$ di velocità e angolo
sia costante lungo l'ultima caratteristica di un ugello ottimale ---
usava un taglio verticale arbitrario attraverso il campo di moto,
aspettandosi che $f_2$ vi variasse. Variava appena: gran parte della
coda regolare di un ugello è quasi piatta in $f_2$, quindi un taglio
verticale è un controllo negativo \emph{debole}. Quello netto è la
\emph{parete}, una linea di corrente lungo la quale $f_2$ sale del
$64\%$; il rigettatore è stato ricostruito su di essa. Lezione: anche
i rigettatori vanno scelti in base alla fisica, non alla comodità.

\section{Passo 4: l'increspatura delle cuspidi di tabella}
\label{app:ripple}

Lo stadio di gradiente del Capitolo~\ref{ch:wallinput} pose dapprima
il test ovvio: confrontare la derivata prodotta dalla
differenziazione automatica (AD) con una differenza finita ---
l'obiettivo ricalcolato dopo aver spostato un ingresso --- a
spostamenti via via più piccoli, aspettandosi che l'accordo
migliorasse. Fallì in un modo che nessuna dimensione di passo poteva
sanare. La differenza finita divergeva man mano che il passo scendeva
sotto la larghezza di una cella delle tabelle di proprietà del gas,
mentre il valore AD oscillava.

Diagnosi: le tabelle sono interpolate linearmente a tratti, quindi la
funzione che il codice implementa davvero ha piccole cuspidi a ogni
confine di cella di tabella, e la derivata \emph{esatta} di quella
funzione porta un'increspatura genuina alla scala delle celle di
tabella --- misurata a una deviazione standard di $44$, cioè lo
$0.010\%$ del gradiente e $1.2\times10^{-8}$ della spinta $J$
(Figura~\ref{fig:ripple}). Nessuno dei due numeri era sbagliato; era
sbagliato il \emph{test}. Un test puntuale AD-contro-differenze
finite su termodinamica tabulata è mal posto.

Il certificato onesto separa la curvatura regolare di $J$ (la fisica)
dall'increspatura (l'interpolazione) adattando una retta su una
spazzata di ampiezze di spostamento e verificando invece l'identità
dell'inviluppo che ne risulta --- il certificato che la suite ora
porta. Una prima analisi aveva anche etichettato erroneamente come
``increspatura'' una curvatura regolare del $2.25\%$; la separazione
per adattamento lineare l'ha corretta al vero $0.010\%$.
Un'interpolazione cubica delle tabelle rimuoverebbe l'increspatura
alla fonte. La stessa trappola attende qualunque solutore
differenziabile costruito su proprietà tabulate, linea dell'aggiunto
CFD compresa.

\section{Passo 7: tre trappole dello strato di ciclo}
\label{app:cycletraps}

Lo strato di ciclo (Capitolo~\ref{ch:cycle}) --- che rappresenta un
RDE con una famiglia di stati di camera stazionari, uno per ogni fase
$\xi$ del passaggio dell'onda --- ha incontrato tre trappole, ora
ciascuna presidiata da un rigettatore.

\begin{enumerate}[itemsep=2pt]
\item \textbf{Lo zero dell'entalpia.} L'entalpia non ha uno zero
  assoluto: ogni tabella fissa il proprio stato di riferimento, e solo
  le differenze sono fisiche. Costruire lo stato di camera di ciascuna
  fase in modo \emph{assoluto}, come $h_0 = h(T_0)$ letto direttamente
  dalla tabella grezza, sbagliava lo stato del $19\%$. Correzione:
  costruire gli stati di fase \emph{relativamente} al riferimento
  proprio delle tabelle
  ($h_0 = h_0^{\rm tab} + h(T_0) - h(T_s^{\rm tab})$, e allo stesso
  modo per l'entropia con il suo termine $R\ln(P_0/p_s)$). Presidio:
  raddoppiare la pressione di camera deve esattamente raddoppiare ogni
  pressione nel campo, $p(2P_0)/p(P_0) = 2$, ora verificato a
  $1.6\times10^{-15}$.
\item \textbf{Il progetto di picco cieco alla pressione.} Il primo
  tentativo di ``progettare alla fase di picco'' teneva fisso il
  numero di Mach di uscita --- il che non è affatto progettare al
  picco, perché ignora dove la pressione di uscita propria di quella
  fase eguaglierebbe l'ambiente. Correzione: il progetto di picco
  risolve $p_e(\varepsilon) = p_a$ per lo stato proprio della fase di
  picco (\code{eps\_star\_at\_pa}).
\item \textbf{Il troncamento silenzioso di tabella.} Una prima
  esecuzione dell'RDE simulato poneva una temperatura di camera di
  picco $T_0 = 4162$~K --- oltre il bordo di 3900~K delle tabelle ---
  e la routine di interpolazione \code{np.interp} \emph{tronca in
  silenzio} al bordo invece di protestare. L'esecuzione si completava,
  con la termodinamica appiattita di nascosto. Correzione: il ciclo
  simulato è stato riposto dentro la validità delle tabelle (picco
  3850~K), e un rigettatore di intervallo di tabella fa ora fallire
  qualunque fase la cui $T_0$ esca dalle tabelle --- la prima cosa da
  controllare quando viene immesso un ciclo misurato.
\end{enumerate}

\section{Passo 7 (assemblaggio): la trappola di segno del getto
libero}
\label{app:signtrap}

Sul bordo di getto libero --- la frontiera dove lo scarico incontra
l'atmosfera, sulla quale la pressione è assegnata e la posizione è
risolta --- la prima versione della cella scriveva la condizione sulla
velocità sostituendo i moduli,
$v = +\sqrt{q_{p_a}^2 - u^2}$ (Figura~\ref{fig:freejet}b). Quella
forma ammette in silenzio solo rotazioni verso l'esterno, perché la
radice quadrata non è mai negativa; un getto in \emph{compressione}
deve ruotare verso l'interno. Il rigettatore del test elementare U-3
--- un bordo in compressione posto su dati esatti --- ha colto la
radice spuria troncata. La cella è stata riformulata con l'\emph{angolo}
di rotazione $\theta_4$ come incognita, così il segno della rotazione
vive dentro l'incognita stessa ed entrambe le direzioni sono
raggiungibili --- la formulazione certificata 6/6 nel
Capitolo~\ref{ch:plug}.

\section{L'assemblaggio del plug: quattro lezioni strutturali}
\label{app:pluglessons}

La marcia del plug del Capitolo~\ref{ch:plug} ha raggiunto la propria
forma certificata attraverso quattro fallimenti strutturali misurati.
Ciascuno è stato colto dall'oracolo in forma chiusa --- una risposta
esattamente nota che la marcia deve riprodurre --- diagnosticato fino
a un meccanismo e chiuso da un cambiamento di progetto. Il progetto
presentato nel Capitolo~\ref{ch:plug} \emph{è} la somma di queste
correzioni.

\begin{figure}[htbp]
\centering
\includegraphics[width=\textwidth]{fig_plug_lessons.pdf}
\caption{Le tre lezioni misurate, disegnate nel mondo di prova
esattamente noto; tutti i pannelli hanno la coordinata assiale $x$
[m] in orizzontale contro la coordinata radiale $y$ [m] in verticale,
con la maglia che la marcia ha effettivamente costruito. (1)~La
radice degenere del bordo: i punti di bordo successivi si accumulano
in un solo luogo invece di avanzare (Sezione~\ref{app:edgeroot}).
(2)~Con un numero fisso di righe di maglia per colonna, la sommità
della maglia si stacca dal getto che si allarga (rosso); aggiungere
una riga per colonna --- lo stesso espediente che usa la marcia della
campana --- mantiene la maglia a riempire il getto (verde). (3)~La
caratteristica che arriva alla parete è un raggio rettilineo dal
labbro: lo stato è costante lungo di esso e varia fortemente
attraverso di esso. Con il piede di quel raggio forzato sulla colonna
adiacente, la relazione di compatibilità veniva chiusa da una singola
regola del punto medio su una corda lunga i cui stati agli estremi
sono interpolati \emph{attraverso} i raggi --- sistematicamente
distorta. La correzione è la ricerca multi-colonna del piede del
Capitolo~\ref{ch:plug}.}
\label{fig:pluglessons}
\end{figure}

\subsection{Lezione 1: la radice degenere del bordo}
\label{app:edgeroot}

Nel primo assemblaggio la cella di bordo libero era alimentata dalla
colonna \emph{precedente}, e nei tratti rettilinei del bordo la
risoluzione di Newton restituiva il nuovo punto esattamente sopra il
vecchio (pt4 = pt3) --- un passo nullo, ripetuto a ogni colonna: il
bordo si accumulava e non avanzava mai.

Il meccanismo (Figura~\ref{fig:edgeroot}) vale la pena enunciarlo
perché è generico. Newton non marcia, \emph{risolve}: qualunque cosa
azzeri i residui è una radice valida, compresa una inutile. In
pt4 = pt3 la condizione di linea di corrente è soddisfatta
identicamente ($0=0$), la condizione di pressione vale già (pt3 sta
sulla frontiera libera per costruzione), e la terza equazione dipende
da che cosa pt3 \emph{sia} nella maglia. Era stato creato, un passo
prima, come uscita di un segmento caratteristico lanciato dai dati di
sommità della colonna precedente --- la maglia registra letteralmente
una caratteristica che termina in pt3. Un mittente preso da quegli
stessi dati giace quindi sulla caratteristica propria di pt3
(esattamente là dove flusso e bordo sono rettilinei, poiché lì i
segmenti discreti giacciono esattamente sulle caratteristiche
rettilinee; al secondo ordine dove si incurvano --- ed è per questo
che il difetto si nascondeva finché il bordo non correva dritto), e
ritracciare quella linea non può che ritrovare la sua uscita
registrata.

La correzione: spazzare ogni colonna dal basso verso l'alto e
alimentare il bordo dal punto interno di sommità della \emph{propria}
colonna --- un punto sulla caratteristica \emph{successiva} della
famiglia, alla cui linea non si chiede affatto di colpire pt3 ma di
attraversare la \emph{linea di corrente} passante per pt3. Quella è
una linea meno inclinata, che una linea più ripida dal basso
intercetta strettamente a valle.

\begin{figure}[htbp]
\centering
\includegraphics[width=0.95\textwidth]{fig_edge_root_geometry.pdf}
\caption{La geometria della radice degenere del bordo; entrambi i
pannelli hanno la coordinata assiale $x$ [m] contro la coordinata
radiale $y$ [m], disegnati alla scala di due celle di maglia.
(a)~L'alimentazione difettosa: pt3 esiste nella maglia solo come
uscita registrata di un segmento caratteristico dai dati di sommità
della colonna precedente, quindi un mittente preso da lì sta sulla
caratteristica propria di pt3 --- ritracciare quella linea non può
che ritrovarne l'uscita registrata, dando pt4 = pt3 e un passo nullo.
(b)~La correzione: il punto interno di sommità della colonna corrente
sta sulla caratteristica \emph{successiva}, e la sua linea deve
attraversare la \emph{linea di corrente} passante per pt3 --- un
bersaglio meno inclinato, intercettato strettamente a valle, così la
radice avanza davvero. Che cosa confrontare: l'angolo fra le due linee
in ciascun pannello.}
\label{fig:edgeroot}
\end{figure}

\subsection{Lezione 2: le righe devono crescere con il getto}

La prima maglia portava un numero fisso di righe per colonna. Ma il
getto si allarga verso valle: le caratteristiche quasi orizzontali di
una famiglia trascinavano in basso le righe fisse mentre il bordo
libero se ne allontanava salendo, finché la sommità della maglia
usciva del tutto dal getto e le due famiglie si aggrovigliavano
(Figura~\ref{fig:pluglessons}, pannello~2). La marcia della campana
sopravvive alla stessa geometria aggiungendo una riga per colonna; la
marcia del plug ora fa lo stesso in modo speculare --- il punto di
bordo libero di ogni colonna diventa una nuova riga di sommità, così
ogni colonna porta una nuova caratteristica e la maglia continua a
riempire il getto.

\subsection{Lezione 3: adiacenza del piede contro accuratezza del
passo}
\label{app:footbias}

Sistemati il bordo e le righe, la cella di parete continuava a
fallire la propria verifica di pressione a valle, e una sonda ha
\emph{nominato} il difetto: angoli di parete esatti a quattro decimali
ovunque, ma velocità di parete distorte e in deriva non appena i
segmenti si allungavano.

Il primo schema confinava il \emph{piede} della cella di parete --- il
punto sulla colonna precedente da cui la caratteristica in arrivo era
partita --- al primo segmento della colonna adiacente. Per tenere lì
il piede le colonne dovevano essere cadenzate, con la spaziatura
assiale libera di crescere fino a un limite di $0.30$. Una singola
chiusura al punto medio su una corda così lunga, che attraversa il
ventaglio di espansione, è sistematicamente distorta: la corda stessa
giace lungo un raggio rettilineo dal labbro, lungo il quale lo stato è
costante, ma i suoi stati agli estremi sono interpolati
\emph{attraverso} i raggi sulla colonna, e la regola del punto medio
propaga quell'errore trasversale sull'intero segmento lungo. I giri
intermedi (colonne allineate, crescita delle righe, cadenza) hanno
portato l'accordo di pressione di parete dal $22\%$ al $55\%$ delle
stazioni e hanno reso il fallimento riproducibile alla stessa
stazione a entrambe le risoluzioni --- la firma di un difetto
sistematico di chiusura anziché di rumore.

La correzione è la ricerca multi-colonna del piede ora nella marcia
(Capitolo~\ref{ch:plug}): le colonne sono collocate alle stazioni
fini assegnate, e la cella di parete cerca sull'intera colonna
precedente, arretrando se necessario, il segmento che racchiude la
caratteristica in arrivo. Le corde si riducono dal limite di $0.30$
alla spaziatura fra stazioni, e l'errore di una chiusura al punto
medio cala con il quadrato della corda. Quel singolo cambiamento ha
portato l'accordo di pressione di parete dal $55\%$ al $98\%$ delle
stazioni e ha riportato in banda l'angolo di bordo.

\subsection{Lezione 4: le righe fantasma, e perché un oracolo solo non
basta}
\label{app:ghostrows}

Con la ricerca del piede a posto la suite segnava ancora 3/5: ogni
cella certificata da Newton, bordo e pressione di parete in banda ---
e la conservazione della massa sbagliata del $+9\%$. La sonda ha
trovato un intero arto di maglia \emph{sotto la spina}, portatore
dello stato uniforme a monte (Figura~\ref{fig:plugbefore}).

Meccanismo: con stazioni fini, la caratteristica di una riga vicina
alla parete termina sulla parete \emph{fra} due stazioni. La marcia
prolungava comunque quella riga, e per una riga simile l'unico
incrocio che Newton riesce ancora a trovare è quello
\emph{all'indietro} --- una radice perfettamente valida delle
equazioni della cella, convergente e certificata, sul ramo sbagliato.
(La ricerca del piede stava già rilevando le terminazioni --- la sua
riga di piede saliva di una riga per colonna --- ma nulla agiva di
conseguenza.) La correzione è il \emph{consumo} delle righe del
Capitolo~\ref{ch:plug}: a una parete solida la famiglia entrante
termina, quindi le righe della colonna precedente sotto il nuovo piede
vengono ritirate e la colonna reindicizzata --- il duale esatto del
bordo che aggiunge una riga per colonna. La massa chiude poi allo
$0.026\%$.

La lezione metodologica è la ragione per cui questa appendice esiste.
Il difetto era invisibile agli oracoli puntuali (P-1, P-2) e alla
certificazione di Newton, perché ogni singola cella era stata
davvero risolta correttamente; solo il bilancio di massa
\emph{globale} lo ha visto. Le suite hanno bisogno di verifiche di
conservazione proprio perché le radici sul ramo sbagliato si
certificano localmente.

\begin{figure}[htbp]
\centering
\includegraphics[width=\textwidth]{fig_plug_before.pdf}
\caption{La regressione dell'arto fantasma, riprodotta dal vivo con
l'interruttore diagnostico del carrier (\code{consume=False}).
(a)~Coordinata assiale $x$ [m] contro coordinata radiale $y$ [m]: le
caratteristiche che avrebbero dovuto terminare alla parete sono
prolungate oltre di essa e fanno crescere un arto di maglia sotto la
spina (punti rossi) --- ogni sua cella certificata da Newton, tutto
quanto sul ramo sbagliato. (b)~La stessa esecuzione misurata sulla
parete: $x$ [m] contro pressione di parete [bar], contro la risposta
in forma chiusa. Le due si sovrappongono --- ed è questo il punto: il
difetto è invisibile alle verifiche puntuali ed è stato colto solo dal
bilancio di massa, sbagliato del $+9\%$. Si confronti con lo stato
certificato, Figura~\ref{fig:plug}.}
\label{fig:plugbefore}
\end{figure}

\section{Passo 10: le lezioni di posizione del ciclo del plug}
\label{app:plugcycleposing}

Lo studio del plug troncato (Capitolo~\ref{ch:plugcycle}) ha raggiunto
la propria forma certificata attraverso tre correzioni misurate ---
due su come il problema debba essere \emph{posto}, e una in cui la
teoria stessa ha rigettato un test trapiantato a sproposito. Tutte e
tre sono disegnate nella Figura~\ref{fig:plugposing}.

\begin{figure}[htbp]
\centering
\includegraphics[width=\textwidth]{fig_plug_posing.pdf}
\caption{Le tre lezioni di posizione del passo 10, tratte dalla
macchina stessa del motore e dai numeri di riferimento.
(a)~Coordinata assiale $x$ [m] contro coordinata radiale $y$ [m]: la
spina a valle. Ponendo lo stesso ventaglio di labbro a
$M_{p_a}=2.35$ il flusso viene ruotato oltre l'assiale
($\theta_E=+15^{\circ}$) e la linea di corrente ``spina'' sale
\emph{sopra} il labbro invece di scendere; la posizione corretta
($M_{p_a}=1.95$, $\theta_E\approx0$) dà una spina genuinamente
discendente. (b)~Fase di ciclo (le cinque fasi, dal passaggio
dell'onda alla coda) contro la pressione di parete che ciascuna fase
produce alla linea iniziale [bar], con l'ambiente condiviso $p_a$
disegnato come linea. A un rapporto di pressione di camera di $2.5$
ogni fase sta \emph{sopra} l'ambiente --- sotto-espansa, e libera di
adattarsi dentro il dominio. Riscalare la stessa scaletta al ciclo
simulato certificato a $PR=10$ porta le fasi basse molto \emph{sotto}
$p_a$: un avvio con urti, che la marcia priva di urti non può
rappresentare. È quel divario a limitare il ciclo utilizzabile.
(c)~La lunghezza di troncamento $l$ [m] come retta numerica, che
segna dove cade ciascuna risposta candidata: l'aspettativa
trapiantata dalla campana (progetto alla fase media, croce rossa) è
smentita dall'ottimo di ciclo misurato $l^{*}=2.103$ --- un
compromesso interno che sta vicino al progetto di picco ma
strettamente al di qua di esso --- con gli ottimi per fase
($\triangledown$) distribuiti nell'ordine che il teorema del plug
prevede.}
\label{fig:plugposing}
\end{figure}

\subsection{La spina a valle}

Il primo mondo poneva la pressione ambiente troppo bassa (la
pressione della fase media a Mach $2.35$). Il ventaglio al labbro del
cowl ruotava allora il flusso di $+40^{\circ}$ --- oltre l'assiale ---
e la linea di corrente destinata a essere la spina si tuffava,
invertiva e saliva \emph{sopra il labbro}: una valle, non una spina,
con la parete che si avvicinava al bordo libero e ogni verifica a
valle che si rompeva sulla geometria. La correzione è la classica
condizione di progetto del plug: porre l'ambiente in modo che la
rotazione del ventaglio cancelli all'incirca l'angolo del flusso
entrante, cioè $\theta_E \approx 0$ (il flusso esce assialmente,
adattato all'ambiente), il che rende la linea di corrente
genuinamente discendente.

\subsection{Il plug sente l'ambiente: il limite di ciclo privo di
urti}

Un ugello a campana non sente mai l'ambiente: il suo flusso è
calcolato fino a un'uscita fissa qualunque sia la pressione esterna, e
$p_a$ entra solo nel punteggio. Un plug lo sente direttamente,
attraverso il bordo libero dove $p = p_a$ è imposta sul flusso
stesso.

Eseguire il ciclo simulato certificato, il cui rapporto di pressione
di camera fra la fase più forte e la più debole è $PR = 10$, contro un
unico ambiente condiviso, metteva le fasi basse $6\times$
sovra-espanse già \emph{alla linea iniziale} --- un regime con urti,
fuori da ciò che una marcia isentropica può rappresentare. Il
fallimento è stato misurato, non ipotizzato: le celle di bordo libero
esplodevano e la spazzatura scendeva riga per riga fino alla parete.
La certificazione di Newton l'ha colto, e la verifica della legge di
scala segnava una dispersione di $14.6$ dove avrebbe dovuto segnare lo
zero di macchina.

Un ciclo per uno studio di plug va dunque posto \emph{privo di urti}:
ogni fase sotto-espansa all'inizio e adattante dentro il dominio. Con
un unico ambiente condiviso questo limita il rapporto di pressione
utilizzabile --- qui $PR = 2.5$. Le affermazioni della teoria non
dipendono da $PR$, e il corpus condiziona già al flusso privo di
urti; ma il limite è un vero limite di modello che vale la pena
dichiarare, perché cicli violenti come $PR = 10$ portano le loro fasi
basse in un funzionamento del plug con urti che questo motore non sa
ancora rappresentare.

\subsection{La teoria che rigetta il test sbagliato, due volte}

Due delle verifiche della prima stesura sono fallite --- ed entrambi i
fallimenti erano la teoria che funzionava, non il codice.

Primo, la stesura trapiantava sul plug il risultato di collasso della
campana (``l'ottimo di ciclo è il progetto per la fase a pressione
media''). È fallito: l'ottimo misurato sta a $l^{*} = 2.10$ contro un
progetto alla fase media di $1.28$. È l'\emph{anti-collasso} --- il
lato opposto della dicotomia della teoria --- perché il risultato di
collasso è dimostrato solo per una parete fissa e interamente
percorsa, che un plug troncato con bordo libero non è. La verifica è
stata invertita nell'oracolo di anti-collasso ora nel
Capitolo~\ref{ch:plugcycle}.

Secondo, la stesura richiedeva che i voti per fase a $l^{*}$ --- la
pressione di parete di ciascuna fase meno l'ambiente sul piano di
troncamento --- variassero monotonamente con la fase e cambiassero
segno una volta, come certificato per la campana nel
Capitolo~\ref{ch:cycle}. Non lo fanno: su un plug le onde di
adattamento fanno oscillare la pressione di parete attorno a $p_a$,
quindi i voti puntuali campionano un'oscillazione e la fase di
passaggio dell'onda può votare negativo a una lunghezza inferiore al
proprio ottimo. La teoria afferma solo il bilancio mediato e la
struttura di compromesso, ed esattamente quelli passano; le verifiche
sovra-specificate sono state sostituite con altre fedeli al teorema
(una dispersione genuina degli ottimi per fase; $l^{*}$ come
compromesso strettamente interno).

La lezione si generalizza: un oracolo può essere trapiantato da una
classe di problemi a un'altra soltanto insieme alle ipotesi che lo
dimostrano.

\section{Passo 15: due strumenti, un vincolo}
\label{app:twoinstruments}

La spazzata sull'angolo di ingresso teneva il proprio vincolo di massa
con una ricerca di radice: spostare il raggio iniziale della spina
finché il flusso che attraversa la linea iniziale non porta la portata
di riferimento. Si valutava poi misurando quella stessa portata. Per
diverse sessioni la valutazione tornava a $7\times10^{-5}$ relativi
--- piccolo, ma un ordine sopra la banda che il raffinamento
sosteneva, quindi la verifica falliva e il fallimento resisteva alle
correzioni ovvie. Allargare la banda l'avrebbe seppellito; raffinare
la quadratura non serviva a nulla.

Il residuo non era affatto un errore. Il vincolo veniva
\emph{imposto} con una regola e \emph{valutato} con un'altra. Il
dimensionamento usava il trapezio del prodotto $\rho\, u\, 2\pi y$; il
verdetto usava il funzionale di colonna proprio della marcia, che
moltiplica medie ai punti medi, $\rho_m u_m\, 2\pi y_m$. Entrambi sono
del secondo ordine, entrambi sono corretti, e differiscono su ogni
intervallo di un quarto del prodotto degli incrementi --- che su $61$
righe è esattamente il $7\times10^{-5}$ che non voleva sparire.

Dimensionare con il funzionale che valuta lo rimuove alla fonte: il
residuo di massa scende fra $10^{-15}$ e $3\times10^{-10}$, cioè il
vincolo è soddisfatto alla precisione di macchina, e alla banda
derivata ($2.9\times10^{-4}$) non si chiede più di coprire un
disaccordo fra due strumenti.

Un secondo fallimento si è dissolto insieme a esso. Un candidato che
ruotava verso l'esterno falliva la certificazione di Newton con un
residuo peggiore di $6.5$, ed era stato letto come un limite di marcia
di quella configurazione. Non lo era: con l'anello dimensionato
correttamente, lo stesso candidato certifica a $0.05$. Un vincolo mal
imposto stava deformando in silenzio la geometria di cui poi si
incolpava la marcia.

La lezione è la forma più affilata di quella del
Capitolo~\ref{ch:compare}: non basta che un vincolo e il suo verdetto
misurino la stessa \emph{grandezza}. Devono essere lo stesso
\emph{funzionale}. Due regole corrette per uno stesso integrale
disaccorderanno al loro ordine di accuratezza comune, e quel
disaccordo si presenterà come un difetto fisico.

\section{Passo 16: tre modi di misurare male con due punti}
\label{app:twopoints}

Costruire la spina a forma libera ha prodotto lo stesso errore tre
volte in una sola sessione, in tre travestimenti. Vale la pena
enunciarlo una volta.

\emph{Primo}, la verifica che la base spline sappia rappresentare il
contorno da cui parte confrontava due conteggi di nodi e leggeva il
loro rapporto come un tasso di convergenza. Il rapporto era
$1.5\times$, troppo lento per un'interpolazione cubica, e la verifica
falliva. Ma sono presenti due errori in competizione --- uno
all'estremo destro libero della spline, uno nella stazione dove il
contorno effettivamente si piega --- e quale dei due domini si alterna
al crescere dei nodi. Due conteggi adiacenti campionano quello che
capita di vincere. Su una scala di quattro conteggi la base riduce il
proprio errore di $28\times$ ed è perfettamente adeguata.

\emph{Secondo}, il vincolo di massa era imposto con una regola di
quadratura e valutato con un'altra
(Appendice~\ref{app:twoinstruments}). Due regole corrette del secondo
ordine per lo stesso integrale disaccordano al loro ordine comune, e
quel disaccordo si presentava come un difetto fisico che nessuna
quantità di raffinamento rimuoveva.

\emph{Terzo}, l'aggiunto era valutato contro una differenza centrata
la cui banda d'errore era stimata dimezzando il passo una volta sola.
Quella stima assume che l'errore sia di troncamento. Qui non lo è: la
marcia riprodotta è solo regolare a tratti nel vettore di progetto,
quindi la differenza porta un rumore che non si riduce con il passo.
Misurata su cinque passi, la discrepanza vaga invece di calare come
$h^2$. La banda a due punti non descriveva dunque nulla, e un aggiunto
corretto veniva riportato come un fallimento.

La struttura comune è questa. \textbf{Due punti possono misurare una
differenza, ma mai un tasso, e mai l'affidabilità dello strumento che
li ha prodotti.} Ovunque questo programma valuti qualcosa contro una
banda derivata, la banda deve venire da una scala --- almeno tre
punti --- e la scala va ispezionata anziché ridotta a un singolo
rapporto. Dove la grandezza sottostante non è regolare, la scala
misura dispersione anziché troncamento, e quella dispersione \emph{è}
la banda onesta.

Vale la pena tenere anche il corollario: tutti e tre i fallimenti
indicavano il colpevole sbagliato. Nel primo si è incolpata la base,
nel secondo il flusso, nel terzo l'aggiunto. In ogni caso la colpa era
dello strumento, e in ogni caso lo strumento era la cosa che non
veniva misurata.

\section{I difetti dell'arbitro, e un binario che non si poteva
ricompilare}
\label{app:geno}

Porre la costruzione classica nel mondo di questo documento
(Sezione~\ref{sec:vsrao}) ha richiesto di eseguire il codice Fortran
indipendente a una massa e a una lunghezza che non gli erano mai state
chieste, e questo ha esposto tre difetti e un fallimento di ambiente.
Sono registrati perché due dei tre erano nostri, e perché il
fallimento di ambiente è del tipo che invalida un arbitro in
silenzio.

\emph{Un NaN dove serviva un rifiuto.} L'integratore della curva di
progetto del codice può esaurire il proprio budget di passi senza
raggiungere uno spigolo, un asse o il proprio bersaglio di massa;
questo è comportamento di base e degradava in silenzio. L'accumulatore
di massa che questo programma ha aggiunto ci passava attraverso e
restituiva una massa NaN. Un chiamante confrontava poi quel NaN con un
bersaglio ottenendo un residuo che non è né positivo né negativo, il
che avvelenava in silenzio una bisezione e alla fine affiorava come
un'entropia NaN nella termodinamica. La correzione verifica gli
incrementi \emph{prima} di accumulare --- sommare e poi sottrarre un
NaN non lo annulla --- e si ferma con una ragione distinta che il
chiamante può rigettare.

\emph{Una soluzione al problema sbagliato.} Con il presidio a posto la
bisezione sulla lunghezza girava, e restituiva un progetto che
rispettava esattamente la lunghezza mancando la massa del $63.6\%$:
aveva trovato la radice su un ramo in cui la curva si ferma al proprio
spigolo anziché al proprio bersaglio di massa. Entrambe le radici
esistono; solo una risolve il problema a due vincoli che era stato
posto. Il test di validità richiede ora la terminazione sulla massa
ogni volta che l'adattamento di massa è attivo.

\emph{Un budget scambiato per un limite.} Il tetto di $500$ passi era
vincolante su curve lunghe. Il progetto che espande fino all'ambiente
di questo documento arriva a $5.8$ m e richiede alcune migliaia di
passi; a $500$ si fermava a metà curva e riportava la lunghezza
raggiunta, ed è per questo che la spina classica per le nostre
condizioni sembrava non esistere. Alzare il tetto non può influenzare
una curva che già termina, e la riproduzione esatta del caso di
riferimento è ciò che lo dimostra.

\emph{Il binario non si poteva ricompilare.} L'eseguibile funzionante
era stato compilato con un compilatore e una libreria di algebra
lineare che non sono più installati. Ricompilare dagli stessi sorgenti
con la toolchain attuale produce un binario che \textbf{fallisce il
caso di riferimento del codice stesso} --- la spazzata sull'angolo di
uscita collassa e il numero di Mach di uscita assegnato diventa
irraggiungibile --- mentre gli stessi sorgenti a $-O0$ riproducono il
riferimento esattamente. Il codice è dunque fragile
all'ottimizzazione con il compilatore attuale, in un modo che il
precedente tollerava.

L'attribuzione conta più della correzione. Il fallimento è comparso
subito dopo le nostre modifiche, e l'inferenza ovvia era che l'avessimo
rotto noi. Ciò che ha risolto la questione è stato rimuovere del tutto
le nostre modifiche e ricompilare: il fallimento si è riprodotto
identico. Solo allora è stato legittimo dire che la colpa era del
compilatore --- e solo allora è stato sicuro tenere la correzione.

\section{Uno strumento che indovinava invece di rifiutare}
\label{app:silentclamp}

Il misuratore di campo del Capitolo~\ref{ch:status} misura massa e
quantità di moto attraverso un taglio verticale, e la sua prima
versione riportava un eccesso di massa del $6$--$7\%$ che decadeva
verso valle. La forma dell'errore era persuasiva: sembrava un difetto
di risoluzione nel campo vicino, e il sospetto ovvio --- il noto buco
di copertura della prima colonna marciata --- è stato provato per
primo. Riempire quel buco non ha cambiato nulla, a nessun
impostazione.

La causa era che il misuratore non sapeva dove fosse il getto. Una
marcia a caratteristiche emette punti di frontiera libera solo una
volta che il fronte di marcia raggiunge la frontiera, quindi la sua
poligonale di frontiera comincia a $x = 2.69$ e non esiste a monte di
lì. La routine che consulta il raggio del getto usava
un'interpolazione standard, che \emph{tronca} in silenzio fuori dal
proprio intervallo di dati e restituiva il primo valore della
poligonale come se fosse il raggio locale. Ogni taglio a monte
copriva allora una fascia esterna fantasma, e poiché l'area va come
$y^{2}$, $0.07$ m di raggio in più a $y \approx 2.3$ valgono qualche
percento della massa.

La diagnosi è stata un singolo numero: in una stazione essenzialmente
sulla linea iniziale, dove la risposta è nota per costruzione a
dodici cifre, il misuratore segnava $+6.35\%$. Uno strumento che
sbaglia dove la risposta è nota sbaglia ovunque.

La riparazione non è un'ipotesi migliore ma un rifiuto: la
consultazione ora solleva un'eccezione fuori dal proprio intervallo.
Ristretto alla finestra in cui la frontiera esiste davvero, il residuo
di massa è dello $0.6$--$1.0\%$ anziché del $7\%$. Un inviluppo
superiore della maglia è stato provato come modo di estendere la
frontiera ed è registrato qui come \emph{sbagliato}: un taglio
verticale attraverso una maglia di caratteristiche è popolato in modo
rado, quindi un massimo per stazione legge i nodi vicini alla parete e
sottostima il raggio del getto fino al $96\%$.
